This work presents a case study on the identification and validation, from real flight data, of the nonlinear dynamic model in multirotor mode of the DH (Hybrid Drone) prototype, a vertical takeoff and landing (VTOL) unmanned aerial vehicle. In this mode takeoff and landing take place, phases that are essential for autonomous operation and demand a representative model for control design. To obtain it, a grey-box formulation is employed, with the structure given by the rigid-body equations and the physical parameters of the aircraft, such as the moments of inertia, the thrust and reaction-torque coefficients and the angular damping coefficients, estimated from the flight data. Since the multirotor is unstable, the data are acquired with the aircraft stabilized by an initial controller, and the plant is identified on its own by relating the input effectively applied by the controller to the measured response, which circumvents the numerical divergence and the correlation between input and states. This process follows the M4V time-domain identification methodology. The excitation maneuvers, of the \textit{doublet} type, are applied to one axis at a time to isolate each channel, and the data undergo synchronization, bias removal and filtering. The estimation takes place in two phases: the equation-error method provides the initial algebraic estimate, refined by the output-error method, which minimizes the difference between the simulated response and the measurements. The identified model is validated on windows distinct from those used for estimation, with compound motions on the three axes, reaching coefficients of determination between 0.70 and 0.77 for the angular rates. The yaw axis proves to be the hardest to identify, owing to its low control authority and high inertia. The model is then linearized around hover, and upon it the cascade control is designed with PD and PI loops, tuned by pole placement. The closed loop is simulated on both the linear and the nonlinear plants, mapping the range of maneuvers over which the linearization remains representative. As its contribution, the work delivers the complete chain from raw flight data to the simulated controller. The validated model surpasses, in its agreement with flight data, the one built solely from geometric measurements and bench tests, and constitutes a reliable basis for vertical takeoff autopilots.

\englishkeywords{System Identification. VTOL. Output-Error Method. Cascade PID Control.}
